\section{Counting flags in the definite row lattices}\label{sec:flag-counts}

An orbital estimate gives the size of one term. Its contribution to the Gram sum also depends on how many integral flag pairs have the given relative position, each counted with its stabilizer weight. We compute
these weights by decomposing the orbits and passing to the
unimodular quotient of the primitive flag. The counts feed
the orbital row sum in Section~\ref{sec:row-sums}. Proposition~\ref{prop:flag-row-mass} applies the genus mass defined below to the marked row; Proposition~\ref{prop:primitive-flag-means} imposes the primitive-line condition on that same mass.

Let \(E/F\) be CM and unramified at finite places, with
\(f=[F:\mathbb Q]\) and \(R=\mathcal O_E\). For an integral Hermitian lattice \((A,h_A)\) define its dual in its span and its discriminant factors by
\[
 A^\#=\{x\in A\otimes_R E:h_A(x,A)\subseteq R\},\qquad \Delta(A)=|A^\#/A|,\qquad d(A)=\sqrt{\Delta(A)}.
\]
For a definite genus \(\mathcal G\) of integral Hermitian lattices its mass is 
\[
 \operatorname{mass}(\mathcal G)=\sum_{[A]\in\mathcal G}|\operatorname{Aut}_R(A,h_A)|^{-1},
\]
 where the sum runs over isometry classes and \(\operatorname{Aut}_R(A,h_A)\) is the group of integral isometries of the Hermitian form on \(A\). Let \(m_j\)
be the mass of the standard unimodular genus of rank \(j\).

\subsection{The row mass and primitive means}

Removing a primitive marked line changes the genus at the primes dividing its norm. We first compute this change for one row, before imposing the conditions of a second flag.

A row \(\Lambda\) is a unimodular lattice with prescribed
intersection \(X_R\) with the isotropic flag span \(X\).
Its middle quotient contains a primitive negative line \(Rz\)
of norm \(-T\), where \(T\geq3\) is completely split. The
off-line definite lattice has rank \(k\). Write \(N=T^f\)
and give each row the reciprocal order of its stabilizer
on the middle quotient.

\begin{proposition}[Mass of the rows]\label{prop:flag-row-mass}
The mass of the possible off-line lattices is
\[
 M=m_k\left(\frac{T^k-1}{T-1}\right)^f.
\]
The total row mass is \((T-1)^f M\). For an invariant
function of the definite lattice, the normalized row
average equals its normalized genus average.
\end{proposition}

\begin{proof}
\proofstep{Splitting off the isotropic flag.}
The primitive submodule \(X_R\oplus Rz\) has projective
quotient. Restriction of the dual and unimodularity give
integral dual partners to \(X_R\) that are orthogonal to
\(z\). Write \(H\) for the Gram matrix of the partners in dual bases.
Subtract \(aH\), regarded as a map back to \(X_R\), for
\(a\in R\) with \(a+\bar a=1\). Surjectivity of unramified
trace at every finite place gives such an \(a\). The new Gram
matrix of the partners is
\[
 H-(aH)^*-aH=(1-\bar a-a)H=0 .
\]
Their pairing with \(X_R\) is intact, so together they span a
unimodular hyperbolic summand.
Carrying it to the fixed hyperbolic summand acts by the
identity on \(X^\perp/X\) and fixes \(z\). Two lifts with
the same middle lattice are therefore equivalent under
an isometry acting trivially on the middle quotient.

\proofstep{Gluing along the discriminant modules.}

The middle lattice is a primitive unimodular extension
of \(Rz\) and its definite orthogonal lattice \(A\) of rank
\(k\). Its gluing group injects in both discriminant modules.
Unimodularity forces the injections to be onto, and
integrality identifies the gluing with the graph of an
anti-isometry of the pairings. Both modules are \(R/(T)\).

At each place over \(T\) the split residue factors of \(R/(T)\)
are paired by conjugation. An isometry multiplies them by \(u,v\)
with \(uv=1\). As there are \(f\) places, the count for a
fixed \(A\) is
\[
    \prod_{\mathfrak t\mid T}|\mathbb F_T^\times|=(T-1)^f .
\]
The single genus of \(A\) has unramified split stabilizer
index \((T^k-1)/(T-1)\) at each place over \(T\). Indeed a
paired split lattice with one elementary divisor of length
one is stabilized by the preimage of a line stabilizer.
Such lattices are obtained by imposing one residue hyperplane
in each pair of factors of a unimodular lattice. Thus, for
\(N=T^f\), the mass is
\[
 M=m_k\left(\frac{T^k-1}{T-1}\right)^f\geq m_kN^{k-1}.
\]
\proofstep{Passing from genus weights to row weights.}

The row weight is the reciprocal order of the stabilizer's
image on the middle quotient. Over \(A\) this is the
stabilizer of the anti-isometry in \(\operatorname{Aut}(A)\).
Let \(\mathcal G(A)\) be the \((T-1)^f\) anti-isometries
over \(A\). For a function \(F\) of its definite lattice, the
contribution of the rows over \(A\) is
\[
 \sum_{[\alpha]\in\operatorname{Aut}(A)\backslash\mathcal G(A)}
        \frac{F(A)}{|\operatorname{Aut}(A)_\alpha|}
   =\frac{|\mathcal G(A)|}{|\operatorname{Aut}(A)|}F(A).
\]
Summing over definite classes gives raw mass \((T-1)^fM\)
when \(F=1\). Dividing the same identity for general \(F\) by
that mass gives the normalized genus average.

\end{proof}

The mass formula determines the average over rows. The pair count also needs the mean number of primitive sublattices of prescribed rank and determinant in this genus.

\begin{proposition}[Primitive sublattice means]\label{prop:primitive-flag-means}
For fixed determinant \(d\), the mean number of primitive
rank-\((k-s)\) sublattices is at most
\[
 (d/N)^k(dN)^{C(s+\log f)}\,\frac{m_{k-s}m_s}{m_k}
 \quad(2\le s\le k-2)
\]
with \(C\) absolute. In ranks one and zero the packing
bound is \((9\mathrm{rd}(E))^{kf}d^k\).
For \(s=1\), \(k\ge3\), and \(d,N\le H\), the mean is at most
\[
 (d/N)^{k-1}(e\pi/(kv))^{kf}h_E
 (4k\sqrt{\mathrm{rd}(E)})^f H^{11+3\log_2 f},
 \qquad v=\mathrm{rd}(E)/2 .
\]
\end{proposition}

\begin{proof}
\proofstep{Counting lattices of fixed index and discriminant.}
Put \(\ell=\log_2(2f)\). A rank \(j\) projective module
has at most \(I^{j+\ell}\) sublattices of index \(I\).
At a prime \(\mathfrak p\) of residue cardinality \(q\), a
sublattice of colength \(a\) is the end of a chain of \(a\)
successive index-\(q\) lattices. At each step the quotient by
\(\mathfrak p\) has \( (q^j-1)/(q-1)\le q^j\) hyperplanes.
Thus the number for any fixed list \((a_{\mathfrak p})\)
of local colengths is at most
\[
  \prod_{\mathfrak p}q^{j a_{\mathfrak p}}=I^j .
\]

For \(p^e\parallel I\), the list above \(p\) satisfies
\(\sum_{\mathfrak p\mid p}f_{\mathfrak p}a_{\mathfrak p}=e\).
Repeat each \(\mathfrak p\) exactly
\(f_{\mathfrak p}a_{\mathfrak p}\) times to encode the list as a
sorted word of length \(e\) in at most \(2f\) letters. This
is injective. Since \(2^e\le p^e\) and \(\ell=\log_2(2f)\),
the number of lists, also bounding the ideal factors of norm \(I\),
is at most
\[
  \prod_{p^e\parallel I}(2f)^e\le
  \prod_{p^e\parallel I}(p^e)^\ell=I^\ell .
\] Maximal integral lattices in a fixed unramified
Hermitian space lie in one genus. At a split place the
two components are dual.

At an inert place diagonalize
using the residual trace, including at two. A maximal
diagonal form has unit entries and at most one uniformizer
entry, since two such entries give a residual isotropic
vector admitting division by the uniformizer. The local
mass correction in the second case is
\((q^j-(-1)^j)/(q+1)\leq q^j\). \cite{gan-hanke-yu}
The mass of a maximal lattice of discriminant order \(D_0\)
is thus at most \(m_jD_0^{j/2}\).
Combine this mass with the inclusion count at index
\(I=\sqrt{D/D_0}\). At the resulting discriminant \(D\) the bound is
\[
     m_jD_0^{j/2}(D/D_0)^{(j+\ell)/2}
         \le m_jD^{(j+\ell)/2}.
\]

\proofstep{Gluing bounds in middle corank.}

For a primitive \(U\subset A\) of rank \(j=k-s\) and
determinant \(d\), let \(B=A\cap U^\perp\). Write \(L=A/(U\oplus B)\) for the gluing group.
It injects in each discriminant module because \(U\) and \(B\)
are primitive in \(A\). The determinant-index formula gives
\[
 \Delta(U)\Delta(B)=|L|^2\Delta(A)=|L|^2N^2 .
\]
Using \(|L|\le\Delta(U)=d^2\) gives
\(\Delta(B)\le d^2N^2\). Witt's theorem gives one orbit per Hermitian type. The determinant-ideal parity at inert places and
positive signatures determine the type by local-global
classification, with at most \((d^2)^\ell\) choices.
\cite{landherr}


Write \(D_U=U^\#/U\) and \(D_B=B^\#/B\). Projection embeds each gluing graph into \(D_B\). At each prime of the integer ring this module is a quotient of a free module of rank \(s\) over a DVR, so its submodules need at most \(s\) generators. Only finitely many primes occur. The Chinese remainder theorem combines their generating lists, padded to length \(s\), into one global list. Padding a generating tuple with zero gives a surjection from a subset of \((D_U\oplus D_B)^s\) to the set \(\mathcal G\) of allowed graphs. Hence
\[
 |\mathcal G|\le(|D_U|\,|D_B|)^s
               \le(d^2\cdot d^2N^2)^s.
\]
An isometry of a glued lattice preserving the marked sublattice restricts to isometries of \(U\) and its orthogonal complement \(B\).

Conversely, a pair of isometries preserves the overlattice precisely when it preserves its graph in the discriminant modules. With \(H=\operatorname{Aut}(U)\times\operatorname{Aut}(B)\), orbit--stabilizer therefore gives the mass at these two isometry classes:
\[
       \sum_{[L]\in H\backslash\mathcal G}\frac1{|H_L|}
           =\frac{|\mathcal G|}{|H|},\qquad H_L=\operatorname{Stab}_H(L).
\]
Apply the two fixed-discriminant mass bounds, sum over
\(\Delta(B)\leq d^2N^2\), and divide by \(M\). This proves
\[
 \operatorname{Mean}\#\{U:d(U)=d\}\leq
 (d/N)^k(dN)^{C(s+\log f)}\frac{m_{k-s}m_s}{m_k}
 \quad(2\leq s\leq k-2)
\]
for an absolute \(C\). In rank one the line lattice
has covolume \(d(\mathrm{rd}(E)/2)^f\).
The product of disks of radius
\[
 R_0=(\mathrm{rd}(E)d^{1/f})^{1/2}
\]
contains a nonzero vector by Minkowski's theorem.
Distinct vectors have
maximum norm separation at least one by integrality and
the norm formula. Packing gives
\[
 (9\mathrm{rd}(E))^{kf}d^k,
\]
which also bounds the rank-zero count.

\proofstep{The corank-one gluing bound.}

For \(s=1\), \(T A^\#\subset A\), and the projection of
\(A^\#\) to the line of \(B\) equals \(B^\#\) by primitivity.
Write \(B_{\mathrm{disc}}=B^\#/B\). The projected gluing
module \(L_B\) is cyclic and satisfies
\[
       T B_{\mathrm{disc}}\subseteq L_B\subseteq B_{\mathrm{disc}}.
\]
There are at most \(N^2\) such choices by counting divisors of
the annihilator of \(B_{\mathrm{disc}}/T B_{\mathrm{disc}}\).
For a generator \(e\) of \(L_B\), evaluation at \(e\) gives the injection
\[
  \operatorname{Hom}(L_B,U^\#/U)\hookrightarrow U^\#/U,\qquad
  \phi\longmapsto\phi(e).
\]
The target has at most \(d^2\) elements. Thus the choices of projected
module and graph map together cost at most \(N^2d^2\).

For one determinant ideal of \(F\) of norm \(d\),
the mass bound in rank \(j\) sharpens to \(m_j d^{j+1}\).
The inert index and mass factor cost together at most
\(q^{ja_v}\) at exponent \(a_v\), and splitting adds
at most \(a_v+1\) choices of the two local indices.
There are at most \(d^{\log_2 f}\) ideals of norm \(d\).
In rank one the inert mass factor is one and the sublattice
of given inert index is unique. Summing line masses up to
determinant \(dN\) costs at most \(m_1(dN)^{2+\log_2 f}\).
These bounds give a numerator at most
\[
 m_{k-1}m_1d^{k-1}d^{4+\log_2 f}N^2(dN)^{2+\log_2 f}.
\]

For every \(k\ge2\), divide the unimodular mass formula
\eqref{eq:row-self-dual-mass} in rank \(k\) by the formula in rank
\(k-1\), with the same definite compact and finite stabilizer
normalizations. Since \(E/F\) is unramified at finite places,
\(|D_E|=|D_F|^2\) and \(|D_F|^{1/f}=2v\). Cancelling the common factors gives
\begin{equation}\label{eq:adjacent-unimodular-genus-masses}
 \frac{m_k}{m_{k-1}}
 =((k-1)!(v/\pi)^k)^f(2v)^{-f/2}L_F(k,\eta^k),
 \qquad v=\mathrm{rd}(E)/2 ,
\end{equation}
with \(\eta\) the quadratic character of \(E/F\).

For \(k\geq3\), \(L_F(k,\eta^k)\geq2^{-f}\), by the
Euler factors, and \((k-1)!\geq(k/e)^k/k\) by integration.
In rank one the unimodular forms in each ideal class
differ by a positive unit modulo norms. Squares are norms,
so \(m_1\leq2^fh_E\). For \(d,N\leq H\) the normalized
corank-one mean is at most
\[
 (d/N)^{k-1}(e\pi/(kv))^{kf}
 h_E(4k\sqrt{\mathrm{rd}(E)})^f H^{11+3\log_2 f}.
\]

\end{proof}

\begin{construction}
\begin{tabular}{@{}lll@{}}\toprule
Gluing datum & Where it is chosen & Number of choices\\\midrule
Cyclic projection \(L_B\) &
\(T B_{\rm disc}\subseteq L_B\subseteq B_{\rm disc}\) & \(\le N^2\)\\
Graph above \(L_B=\langle e\rangle\) &
\(\phi\mapsto\phi(e)\in U^\#/U\) & \(\le d^2\)\\\bottomrule
\end{tabular}
\caption{The two choices in corank-one gluing. Each projected module admits at most \(d^2\) graph maps, so the product bound is \(N^2d^2\).}\label{fig:corank-one-gluing-data}
\end{construction}

\subsection{Integral frame classes over a sublattice}

The cross pairings fix a rational relative position. To count integral representatives of that position, one must bound the choices of frames that give it over a prescribed sublattice.

Write \(Z\) for the cross Gram of the augmented frames,
\(\Delta_j\) for its leading minors with \(\Delta_0=1\), and \(Y\) for the
first \(r\) rows. The relevant cells express \(Z\), up to
column signs, as unitriangular factors around a block
anti-diagonal monomial matrix with constant scalar on
each reversed block.

\begin{lemma}[Heights of the reduced frame]\label{lem:reduced-flag-heights}
Write \(Z=LWU\), with unitriangular \(L,U\) and a scalar times reversal
on each diagonal block of \(W\). Assume nonzero Gram invariant,
\(\mathrm{rd}(E)\le b\), \(T,T'\le b^{C_Tm}\), and \(f\le b^{C_g}\).
The first \(r\) rows have a nonzero full-row minor \(\delta\).
The reduced frame has the direct and reciprocal normalized
height bounds \(H_C(Z)\) of \eqref{eq:reduced-flag-height}
for a sufficiently large fixed \(C\),
and the three arithmetic factors satisfy 
\[
 \max\!\left\{\operatorname N_E(\delta),\ \prod_\tau|\tau G|,\ \operatorname N_E(\mathfrak d)^r\right\}\le H_C(Z)^2.
\]

\end{lemma}

\begin{proof}
Deleting the last row and its
monomial pivot column gives a full row minor \(\delta\).
For the last block, of size \(r+1-l\) and pivot scalar \(t\),
the two determinant identities are
\[
 |\delta|=|\Delta_l||t|^{r-l},\qquad
 |\Delta_{r+1}|=|\Delta_l||t|^{r+1-l}.
\] In the leading \(l\) positions all pivots are retained.
At the last block, take precisely the pivot columns of the
remaining \(r-l\) rows. In the Cauchy--Binet expansion, the
triangular factors force this set of intermediate columns.
Their principal triangular minors have determinant one. Thus
the selected determinant has exactly \(r-l\) factors of \(t\)
after the leading \(\Delta_l\), while the full determinant
has \(r+1-l\).

Integral operations in decreasing pivot order reduce
both triangular factors to entries of size at most
\(\sqrt{2f}\,\mathrm{rd}(E)\). This covering bound follows
from Euclidean transference and the inverse different
description of the dual. \cite{banaszczyk}
Cauchy--Binet bounds every minor of the reduced
representative by
\[
 4^{r+1}((r+1)!)^2(\sqrt{2f}\,\mathrm{rd}(E))^{2(r+1)}
 \prod_{j=1}^{r+1}\max(1,|\Delta_j|)\max(1,|\Delta_j|^{-1})
\]
at each embedding, with a zero minor contributing one.
The block factors are bounded by the displayed product
as each block gives the ratio of its two boundary
determinants. The product formula bounds reciprocal
products for all nonzero integral minors and for
\(G=|\Delta_{r+1}|^2-TT'|\Delta_r|^2\) by the direct
products. Here \(|x|^2=x\bar x\). The frame weights are
at most one with reciprocal products at most
\(\mathrm{rd}(E)^{2(r+1)^2}\).

At the embedding \(\tau\), let \(G_{z,\tau}\) be the scaled Gram determinant and let \(d_{j,\tau}(Z)\) be the absolute value of the scaled leading minor of order \(j\).
For \(m=2r+2\) set
\begin{equation}\label{eq:reduced-flag-height}
 H_C(Z)=b^{C^2m^3f}\left(\prod_{\tau,j}
 \max(1,d_{j,\tau}(Z)/b^{Cm^2})\right)^4.
\end{equation}
Under \(\mathrm{rd}(E)\leq b\), \(T,T'\leq b^{C_Tm}\),
and \(f\leq b^{C_g}\), these direct and reciprocal
normalized products are at most \(H_C\), for fixed
sufficiently large \(C\). The unnormalizing powers and
cutoff cost at most \(b^{(C+1)m^3 f+O_{C_g,C_T}(m^2f)}\).
For \(Q=H_C^2\) we also have
\(\operatorname N_E(\delta),\prod_\tau|\tau G|,
\operatorname N_E(\mathfrak d)^r\leq Q\), where
\(\mathfrak d\) is the different.

\end{proof}

\begin{proposition}[Frames over a fixed sublattice]\label{prop:integral-flag-frames}
In the nonzero-Gram case, assume the height bounds of
Lemma~\ref{lem:reduced-flag-heights} and put \(Q=H_C(Z)^2\). If \(P\subset A\) has rank
\(k-s\) and determinant at most \(D\), then
\[
 \#\mathrm{Fr}(Z,P)
 \le N^{2r}D^r\,9^{sf}(D^2N^2Q)^{s+sr}Q^{2r^2}.
\]
Here \(\mathrm{Fr}(Z,P)\) consists of terms of Gram class \(Z\)
whose orthogonal lattice contains \(P\) after projection along
the isotropic span.
\end{proposition}

\begin{proof}
\proofstep{Reducing to a specified integral lift.}
Choose a split row and an exact Gram representative.
The right integral upper unitriangular operation
reparametrizes the second frame. The adjoint of the
left operation preserves the inverse different powers
on \(X_R\), and extends by inverse adjoint on the dual
partners and a shift of \(z\) into \(X_R\).
For an integral map \(S\) from the middle lattice to \(X_R\), use the two transformations
\[
    w\longmapsto w+Sw,\qquad
    y\longmapsto y-S^*y-aSS^*y .
\]
The first sends middle vectors into the graph of \(S\).
Its new cross pairing with the dual partners vanishes because
of \(-S^*\). The new partners' Gram vanishes by \(a+\bar a=1\),
so this is the required integral isometry.

\proofstep{Counting lifts and cofactor vectors.}

The lift \(p+\phi(p)\) of \(p\in P\) is unique because
\(Y\) has full row rank. Shifts of \(\phi\) come from
maps \(A^\#\to X_R\), as the middle lattice projects
onto \(A^\#\). In local invariant-factor coordinates for \(P\subset A^\#\cap EP\),
restriction to \(P\) multiplies each lift coordinate by the associated
elementary divisor. There are \(r\) target copies, so the number of
lift classes is at most
\[
        [A^\#\cap EP:P]^r\le N^{2r}D^r .
\] Fix one lift. Let \(S_0=A\cap P^\perp\), of rank
\(s\), with \(\Delta(S_0)\leq D^2N^2\) by the gluing
bound. Its discriminant module has annihilator \(I\)
of norm at most \(D^2N^2\). The frame projections
\(w_b\) to \(P^\perp\) lie in \(S_0^\#\).

For the signed cofactor vector \(c\) of \(Y\),
\[
 T h\!\left(\sum c_bw_b,\sum c_bw_b\right)=G .
\]
Indeed the same combination of the frame vectors lies
in \(X^\perp\), has norm \(-T'|\Delta_r|^2\), and pairs
with \(z\) by a signed \(\Delta_{r+1}\).
Orthogonality to the lift puts its positive part in
\(P^\perp\). Thus \(G>0\) at every embedding in the
nonzero-Gram case. For a difference \(x\ne0\) of cofactor combinations,
\(\bar I h(x,x)\subset R\). As the norm of the nonzero integral
ideal \(h(x,x)\bar I\) is at least one, the product over the
chosen complex embeddings satisfies
\[
 \prod_\tau h_\tau(x,x)\ge (\operatorname N_E I)^{-1/2}.
\]
Divide the squared norms by the squared radii \(|\tau G|\).
Their geometric mean gives normalized separation
\((D^2N^2Q)^{-1/(2f)}\). Packing in the resulting product of
balls gives at most \(9^{sf}(D^2N^2Q)^s\) possibilities.

\proofstep{Counting the remaining frame entries.}

Fixing the integral lift removes the ambiguity from the common sublattice. The residual multiplicity is now a count of entries of the two frames with specified pairings.


At fixed cofactor combination tuple differences are
\(vY\), where Cramer's rule puts \(v\) in
\((\delta^{-1}S_0^\#)^r\). Integral Eichler shifts
from \(S_0^r\) realize changes by \(vY\) and fix the
lift. There remain at most \((D^2N^2Q)^{sr}\)
classes. At fixed projections, two frames differ
only in \(X\). For the difference matrix \(B\), the pairing identity gives
\[
 B^*Y+Y^*B=0\quad\Longrightarrow\quad Y^*Bz=0\qquad(z\in\ker Y).
\]
 Since \(Y\)
is onto, \(Y^*\) is injective. Thus \(B=CY\), and
\[
       0=Y^*(C^*+C)Y\quad\Longrightarrow\quad C^*=-C.
\] The entries of \(C\) lie in
\(\delta^{-1}\mathfrak d^{-r}\) by Cramer's rule.
Integral skew matrices act by central isotropic
shifts, leaving at most \(Q^{2r^2}\) classes. Thus
\begin{equation}\label{eq:flag-over-P}
 \#\mathrm{Fr}(Z,P)\leq
 N^{2r}D^r9^{sf}(D^2N^2Q)^{s+sr}Q^{2r^2}.
\end{equation}

\end{proof}

\subsection{Pairs and the generic determinant}

A pair has more structure than either frame separately. Their common sublattice bounds the ambiguity in assembling them, while its orthogonal determinant distinguishes the generic count.

\begin{proposition}[Pair counts through an integral sublattice]\label{prop:flag-pair-counts}
For the paired terms with full row rank in the projection and
\(d(M_0)I_1I_2\le D\), the common sublattice has discriminant
size at most \(D\) and the two term classes determine the pair
class under the stated congruence conditions. Write \(H=H_XH_Y\)
for the product of the reduced label heights and \(N=T^f\). For
\(2\le s\le k\), \(D\le H^3\), the count per unit of total row mass is at most
\begin{equation}\label{eq:flag-pair-mass-factor}
 (D/N)^k H^{E_s}(csf)^{6sf}\Xi_s(H),\qquad
 E_s=c\big((s+m)m+2+\log_2 f\big),
\end{equation}
where \(c\ge1\) is an absolute counting constant and
\[
 \Xi_s(H)=
 \begin{cases}
    m_{k-s}m_s/m_k,&2\le s\le k-2,\\
    (ck)^{kf}(HN)^{ck},&s=k-1,k .
 \end{cases}
\]
In the generic
case the orthogonal determinant is

\[
 D_0=g^{-mf}\prod_\tau G_{z,\tau},
\]
 and at common rank \(k-1\)
the full orthogonal lattice of the first term can be used in the
partner bound.
\end{proposition}

\begin{proof}
\proofstep{Finding the common integral sublattice.}

For two terms the projected ordinary orthogonal lattices
\(U_i\) lie in the projected enlarged congruence lattices
\(A_i\) with the same ranks and indices \(I_i\) as before
projection. The projection is injective on the rational
orthogonal complement because \(Y\) has full row rank.
The pair correspondence gives \(M_0\subset A_1\cap A_2\).
Indeed the paired vectors differ by \(X\), and the
comparison isometry transports the second representative
of the level relation while preserving its congruence.
The map \(x\mapsto(x+U_1,x+U_2)\) now gives an injection
\[
 M_0/(M_0\cap U_1\cap U_2)\longrightarrow A_1/U_1\oplus A_2/U_2 .
\]
With \(P=M_0\cap U_1\cap U_2\), the injection bounds
\([M_0:P]\) by \(I_1I_2\). In the determinant convention here,
\[
       d(P)=d(M_0)[M_0:P]\le d(M_0)I_1I_2\le D.
\]

\proofstep{Eliminating the comparison isometry.}

The pair class is determined by its term classes.
The relative comparison isometry of a fixed pair
acts trivially modulo \(g\) on the definite quotient
of an augmented flag complement by its isotropic span.
For an integral \(x\) in the complement,
the level congruence of the two isotropic lattices
makes \(h(x,X_R)\) divisible by \(g\). Unimodularity
extends the integral pairing \(h(x,X_R)/g\) to a vector \(l_0\)
of the row lattice \(\Lambda\). Therefore
\[
       l=x-g l_0\equiv x\pmod{g\Lambda},\qquad h(l,X_R)=0 .
\]

Applying
pair compatibility to \(l,x\) under both maps gives
the congruence on the quotient. The quotient isometry
has finite order by definiteness. The integral
principal congruence subgroup at \(g\geq3\) is
torsion free, by the binomial expansion at an odd
prime dividing \(g\), or modulo \(4\) for a power of
two. The maps also agree on the negative line modulo
the isotropic space, which proves uniqueness on the
entire quotient used in the pair class.

Multiply \eqref{eq:flag-over-P} for the two terms.
\proofstep{Summing over saturations.}

The common lattice need not be saturated. Summing over its possible saturations is what turns the discriminant bound into a bound for integral pairs.


For a saturation of determinant \(d\), at most
\((D/d)^{k+\ell+1}\) sublattices lie over it.
In the primitive means, multiplication by the sublattice bound gives
\[
     (d/N)^k(D/d)^{k+\ell+1}=(D/N)^k(D/d)^{\ell+1}.
\]
If
\(H=H_C(Z')H_C(Z'')\) and \(D\leq H^{A_d}\)
with fixed \(A_d\), use \(Q\leq H^2\) in the two
counts and absorb all residual powers in
\(H^{C((s+m)m+1+\ell)}(Csf)^{6sf}\).
Use the middle-corank mass comparison when \(s\le k-2\) and
the packing bound in the other range. The extra factor is thus
\[
   \begin{cases}
     m_{k-s}m_s/m_k,&s\le k-2,\\
     (Ck)^{kf}(HN)^{Ck},&s>k-2.
   \end{cases}
\]

\proofstep{Computing the generic determinant.}

The generic term has an additional determinant saving.
The Gram of the two augmented frames has determinant
\((-1)^{r+1}G\) by the rank-one block determinant
identity. Incorporating both flag lattices' different
powers gives determinant \(\prod_\tau G_{z,\tau}\).
The primitive hull quotient is zero away from \(g\).
At primes over \(g\), genericity generates it by the
classes of \(g^{-1}(e'_b-\sum e_a\xi_{ab})\), with
augmented frames. These generators are independent
over \(R/(g)\) because the rational flag spaces are
complementary and the different powers are units
at those primes. The index is \(\operatorname N_E(g)^{r+1}=g^{mf}\).

By unimodular gluing the hull and its orthogonal complement have
equal discriminant modules in size. In the determinant convention
\(d\), division by the hull index gives
\[
       D_0=g^{-mf}\prod_\tau G_{z,\tau}.
\]
Use \eqref{eq:flag-over-P} with \(s=1\), \(D=D_0\)
and the corank-one mean. The same divisor sum gives
\((D_0/N)^k\) at the cost of one factor \(N/d\).
The residual powers are at most \(H_C^{Cr^2}\).
Here \(h_E\leq\mathrm{rd}(E)^{f(1+\ell)}\) by Minkowski
and ideal counting, and \(\log f\leq\sqrt r\) and
\((kb)^8\leq g\leq b^{C_g}\) absorb the class factor.
For a reduced label \(Z\), the normalized mean number of generic
first terms is therefore bounded by \(H_C(Z)^{Cr^2} n_Z\), where
\begin{equation}\label{eq:generic-first-term-factor}
 (e\pi/(kv))^{kf}(D_0/N)^k
 =\prod_\tau\left(\frac{e\pi G_{z,\tau}g^{-m}}{Tkv}\right)^k=n_Z .
\end{equation}

\proofstep{Recovering integrality of the partner.}

For a fixed first generic term a partner of common
rank \(k-1\) contains its entire orthogonal lattice
modulo \(X\). Indeed every integral \(x\) of the first
complement has a rational partner \(y\).
Compatibility writes \(x-\beta x=g l+x_0+x_2\), where
\(l\in\Lambda\) and the \(x_i\) are integral isotropic vectors
of the indicated flags. The two rational flag components of
\(g l\) are
\[
     g l=(x-y-x_0)+(y-\beta x-x_2).
\]
The vector \(l\) lies in the primitive hull of the second term's
two flag lattices. Genericity puts \(g l\) in their integral sum.
Its first component \(x-y-x_0\) is consequently integral, so
\(y\) is integral too.
The fixed lattice has determinant at most \(H_C(Z')\),
so \eqref{eq:flag-over-P} with \(s=1\) bounds the
partner count by \(C^f(H_C(Z')H_C(Z''))^{Cr^2}\).

\end{proof}

\subsection{Covolume of the simultaneous flag stabilizer}

The multiplicity count leaves one measure factor in the unfolding. Its size is controlled by writing the common stabilizer in coordinates adapted to the same pair of flags.

\begin{proposition}[Covolume in root coordinates]\label{prop:flag-stabilizer-covolume}
Assume that \(Z\) has nonzero Gram invariant and the block
factorization and height bounds of Lemma~\ref{lem:reduced-flag-heights}. Write \(R_{\rm stab}\) for the stabilizer lattice of the two flags in the term normalization. For a sufficiently large fixed \(C\), 
\[
 \operatorname{covol}(R_{\rm stab})\le H_C(Z)^{Cr^2}.
\]

\end{proposition}

\begin{proof}
Put \(n=r+1\) and write the sign-corrected Gram as \(Z=LWU\).
For an upper unitriangular \(a\), the other triangularity condition
becomes
\[
 ZaZ^{-1}\text{ lower unitriangular}
 \quad\Longleftrightarrow\quad
 W(UaU^{-1})W^{-1}\text{ lower unitriangular}.
\]
Conjugation by \(L\) preserves the lower unitriangular group.
Conjugation by \(W\) then selects exactly the upper roots within
its reversed blocks. These root coordinates
have Lebesgue Haar measure, since multiplication in
order of increasing root height is triangular with
diagonal entries one. The column signs preserve it.

\proofstep{Integral lifts of root entries.}

If the entries of \(a-1\) lie in \(I=(G^2\Delta_n)\),
the pair \(b,a\) acts on the two frames, with
\[
      b=(Za^{-1}Z^{-1})^*,\qquad b^*Za=Z .
\]
The upper-triangular actions with fixed last diagonal entry preserve
the internal forms, and the matrix identity preserves the cross
Gram. Extend by identity on their
common orthogonal complement. Let \(A_{\rm joint}\) be the Gram matrix on the common frame
span. In its projection the only possible denominator divides
\(\det A_{\rm joint}=(-1)^nG\). Cramer's rule gives
\[
  a-1, b-1, a^{-1}-1, b^{-1}-1\in G M_n(R).
\]
The changes on the two frames therefore clear the denominator
of the projection. The resulting isometry preserves the
integral row. It acts unipotently on the isotropic
flags, and trivially on the quotient by each
isotropic span in its perpendicular space.
Indeed it fixes the negative line there and is
unipotent and unitary on the positive complement.

\proofstep{Denominators in the triangular coordinates.}

Let \(d_i\) be the nonzero leading determinants at
the block boundaries, \(d_0=1\), and \(P_0=\prod_{i>0}d_i\).
Block Gaussian elimination gives a factorization with
identity diagonal blocks in \(L\). In block row \(i\),
the upper factor is \(W_i^{-1}\) times the Schur row,
whose entries have denominators dividing \(d_{i-1}\).
For \(j>i\) the off-diagonal block in column \(i\) is the
last block of
\[
   Z_{j,\le i}Z_{\le i,\le i}^{-1}
     =d_i^{-1} Z_{j,\le i}\operatorname{adj}(Z_{\le i,\le i}).
\]
Its denominator thus divides \(d_i\).

For the scalar \(t_i\) of a block of size \(h_i\), use
\[
 (d_{i-1}t_i)^{h_i}=\pm d_{i-1}^{h_i-1}d_i,\qquad
 (d_i/t_i)^{h_i}=\pm d_i^{h_i-1}d_{i-1}.
\]
The elements \(d_{i-1}t_i\) and \(d_i/t_i\) are integral by
integral closedness. Thus \(P_0^2U\) and \(P_0^2LW\) are integral,
and \(U^{-1}=Z^{-1}LW\) makes \(\Delta_nP_0^2U^{-1}\) integral. Any other right triangular factor
changes the root coordinates by an inner conjugation
and preserves their measure.

\proofstep{Fundamental domains and Haar volume.}

The bounded denominators determine a root lattice. Its volume in the prescribed measure is the factor needed in the orbital integral.


The level ideal \(J=I(\Delta_nP_0^4)\) on the root
entries thus gives integral simultaneous lifts.
Use root translations in increasing height to put
each entry in an additive fundamental domain of
the embedded \(J\). Each additive fundamental domain has covolume
\((\mathrm{rd}(E)/2)^f\operatorname N_EJ\). The product has volume
\[
       \big((\mathrm{rd}(E)/2)^f\operatorname N_EJ\big)^{d_W},
       \qquad d_W\le n^2 ,
\]
where \(d_W\) counts the roots.
Writing \(S=\prod_{\tau,j}\max(1,|\Delta_j|_\tau)\),
we have \(\operatorname N_EJ\leq(2TT')^{4f}S^{20}\).
After multiplying by the measure factor
\(\mathrm{rd}(E)^{2fn^3}\) of the term normalization,
this volume is at most \(H_C(Z)^{Cr^2}\) for fixed
sufficiently large \(C\), by the same size conversion
as for the Gram minors.
\end{proof}
